\documentclass[11pt]{article}
\usepackage[margin=1in]{geometry}
\usepackage{amsmath,amssymb}
\usepackage{booktabs}
\usepackage{hyperref}
\usepackage{xcolor}
\usepackage{graphicx}

\title{Independent Replication Report:\\
\emph{Quantum skyrmions in frustrated ferromagnets}\\
\large Lohani, Hickey, Masell \& Rosch, \textit{Phys.\ Rev.\ X} \textbf{9}, 041063 (2019); arXiv:1901.03343}
\author{REPLICATE-PROJECT / TEXTURES-100 --- independent from-scratch exact diagonalization}
\date{}

\newcommand{\Svec}{\mathbf{S}}
\newcommand{\Cperp}{C_{\perp}}

\begin{document}
\maketitle

\begin{abstract}
We independently reproduce the central exact-diagonalization (ED) result of Lohani \emph{et al.}
(2019): that a quantum skyrmion in a frustrated spin-1/2 XXZ triangular-lattice ferromagnet
exists as a \emph{stable many-magnon bound state}. Working from the paper's equations only (no
author code), we build the many-body Hamiltonian by hand in fixed-magnetization ($S_z$) sectors
and diagonalize with \texttt{scipy.sparse.linalg.eigsh}. We confirm (i) a negative binding energy
$E_0^B(N_f)<0$ for every $N_f\ge2$ on both the 7-site and 19-site round flakes --- the rigorous
finite-size signature of magnon binding --- and (ii) the quantitative anchor: the maximal
transverse antiferromagnetic correlation reaches $\Cperp = 0.73$ (19-site flake, $N_f=4$), landing
squarely inside the paper's quoted $\Cperp \approx 0.6$--$0.8$ window for proper quantum skyrmions
(Figs.~2,~5). We further document that the raw scalar chirality
$\langle\chi\rangle=\sum_\triangle\langle\Svec_i\cdot(\Svec_j\times\Svec_k)\rangle$ is
\emph{exactly zero} --- and explain why this is correct physics, not a replication failure.
Finally, in a finite-size-scaling extension we rerun the ED on progressively larger flakes
($N=7,19,37$; $37$ approaches the paper's 31-site cluster) with low-lying spectra
($k>1$ via Lanczos) and reproduce the paper's \emph{headline dynamical claim}: the
skyrmion$\leftrightarrow$antiskyrmion \emph{tunneling splitting} $\Delta_{\rm tun}=E_1-E_0$ within
the skyrmion sector is \emph{exponentially small}, collapsing from a generic spacing $\sim0.23$
(no skyrmion, $N=7$) to $\sim10^{-2}$ and down to machine-zero $\sim10^{-14}$ ($N=19,37$), with the
predicted $N_f\bmod3$ selection structure.
\textbf{Verdict: REPLICATED} (core ED claim + quantitative anchor + tunneling-bandwidth FSS),
Coverage $\approx 7/10$, Agreement $\approx 8/10$.
\end{abstract}

\section{The paper and its central claim}
The paper develops a \emph{quantum} theory of magnetic skyrmions and antiskyrmions in a spin-1/2
Heisenberg magnet whose ferromagnetism is frustrated by antiferromagnetic next-nearest-neighbor
(NNN) exchange. Its model (Eq.~1) is the XXZ triangular-lattice Hamiltonian in a field:
\begin{equation}
H = -J_1 \sum_{\langle ij\rangle} \Svec_i\cdot\Svec_j
    + J_2 \sum_{\langle\langle ij\rangle\rangle} \Svec_i\cdot\Svec_j
    - K \sum_{\langle ij\rangle} S_i^z S_j^z
    - B \sum_i S_i^z ,
\label{eq:H}
\end{equation}
with $J_1$ a ferromagnetic (FM) nearest-neighbor (NN) coupling set to $1$, $J_2$ an
antiferromagnetic NNN coupling (frustration), $K>0$ an easy-axis anisotropy that stabilizes
skyrmions, and $B$ an external field along $z$. The model is spin-rotation invariant about $z$.

Because a topological winding number cannot be sharply quantized in a quantum spin system
(Heisenberg uncertainty), the authors \emph{define} a quantum skyrmion operationally: a
\emph{stable many-magnon bound state} whose correlation functions smoothly connect to the classical
skyrmion texture. The headline ED result is therefore not a global ground state (the polarized FM
wins globally at these fields) but a \emph{bound state within a fixed-magnetization sector},
stable against ``quantum evaporation'' of magnons.

\paragraph{The two checkable quantities.}
\begin{enumerate}
\item \textbf{Binding energy.} Let $N_f$ be the number of flipped spins (magnons) relative to the
fully polarized FM, i.e.\ $N_f = S_z^{\mathrm{FM}} - S_z$. Define the single-magnon energy
$e_1 = E(N_f{=}1) - E_{\mathrm{FM}}$ and the binding energy
\begin{equation}
E_0^B(N_f) = E_0(N_f) - N_f\, e_1 .
\label{eq:binding}
\end{equation}
$E_0^B(N_f)<0$ means the $N_f$ magnons prefer to bind into a single droplet rather than disperse
as $N_f$ independent magnons --- this \emph{is} the quantum skyrmion. Crucially $E_0^B$ is
\emph{field-independent} (the Zeeman term shifts every state in a sector uniformly), so we compute
it at $B=0$. A negative $E_0^B$ on a finite flake is a rigorous variational upper bound and hence
proves binding on the infinite lattice.
\item \textbf{Transverse anticorrelation.} The skyrmion texture is detected by the maximal
transverse antiferromagnetic correlation
\begin{equation}
\Cperp = \max_{ij}\; -\big\langle S_i^+ S_j^- + S_i^- S_j^+\big\rangle
       = \max_{ij}\; -2\big\langle S_i^x S_j^x + S_i^y S_j^y\big\rangle .
\label{eq:cperp}
\end{equation}
The paper quotes $\Cperp \approx 0.6$--$0.8$ for proper quantum skyrmions (Figs.~2,~5). This is the
one clean quantitative number to hit.
\end{enumerate}

\section{Method: from-scratch exact diagonalization}
We reimplement the model from Eq.~\eqref{eq:H} in Python (\texttt{numpy}/\texttt{scipy}); no author
code was used or available (none is published). The implementation
(\texttt{report/evidence/lohani\_ed.py}) proceeds as follows.

\paragraph{Geometry.} Approximately round triangular flakes are built from axial coordinates: all
lattice points within radius $R$ of the origin. $R=1$ gives the $N=7$ site flake (hexagon $+$
center; 12 NN bonds, 6 NNN bonds, 6 elementary triangles); $R=2$ gives the $N=19$ site flake (42 NN,
30 NNN, 24 triangles). NN bonds sit at distance $1$, NNN at $\sqrt{3}$; elementary triangles are the
shared-neighbor pairs of each NN bond.

\paragraph{Many-body basis.} We exploit the conserved total $S_z$: for a sector with $n_\downarrow$
down spins, the basis is all $\binom{N}{n_\downarrow}$ bitstrings with exactly $n_\downarrow$ set
bits (bit $=1 \Rightarrow$ down $\Rightarrow S^z=-\tfrac12$). A dictionary \texttt{\{state:index\}}
provides the column lookup for off-diagonal matrix elements.

\paragraph{Hamiltonian assembly.} Using
$\Svec_i\cdot\Svec_j = S_i^z S_j^z + \tfrac12(S_i^+S_j^- + S_i^-S_j^+)$, the diagonal carries the
$zz$ exchange, the anisotropy, and the field; the off-diagonal flip-flop $S_i^+S_j^-$ connects
states differing by one up$\leftrightarrow$down swap with amplitude $\tfrac12\times$coupling. The
sector Hamiltonian is assembled as a \texttt{scipy.sparse.csr\_matrix} and its lowest eigenpair
found via \texttt{eigsh(k=1, which='SA')} (dense \texttt{eigh} fallback for tiny sectors). Sweeping
$n_\downarrow = 0\ldots N/2$ and minimizing over sectors gives the global ground state; the binding
analysis instead compares each fixed-$N_f$ sector.

\section{Results}

\subsection{Binding energy: many-magnon bound state confirmed}
Table~\ref{tab:binding19} lists the 19-site flake ($J_2=0.5$, $K=0.05$) per-sector energies.
$E_0^B(N_f)<0$ for \emph{every} $N_f\ge2$ and deepens monotonically with $N_f$ --- the magnons bind
into a droplet exactly as Eq.~\eqref{eq:binding} demands. The excitation energy above the FM,
$\Delta E = E_0(N_f)-E_{\mathrm{FM}}$, even turns \emph{negative} at $N_f=5,6$
($\Delta E=-0.031,-0.096$), i.e.\ the multi-magnon droplet drops below the FM in that regime,
matching the paper's picture of a droplet forming with $N_f^{\min}\sim5$--$9$. The 7-site flake
shows the same qualitative pattern with a minimum binding of $-0.60$ ($J_2=0.5$) and $-0.70$
($J_2=0.7$).

\begin{table}[h]
\centering
\begin{tabular}{rrrrr}
\toprule
$N_f$ & $E_0(N_f)$ & $\Delta E$ & $E_0^B(N_f)$ & $\Cperp$ \\
\midrule
0 & $-7.2750$ & $ 0.0000$ & $ 0.0000$ & $0.00$ \\
1 & $-7.2549$ & $ 0.0201$ & $\approx 0$ & $0.00$ \\
2 & $-7.2483$ & $ 0.0267$ & $-0.0134$ & $0.48$ \\
3 & $-7.2567$ & $ 0.0183$ & $-0.0419$ & $0.60$ \\
4 & $-7.2726$ & $ 0.0024$ & $-0.0779$ & $\mathbf{0.73}$ \\
5 & $-7.3058$ & $-0.0308$ & $-0.1312$ & $0.63$ \\
6 & $-7.3715$ & $-0.0965$ & $-0.2170$ & $0.32$ \\
\bottomrule
\end{tabular}
\caption{19-site round flake, $J_1=1$, $J_2=0.5$, $K=0.05$, $B=0$. $e_1=0.02009$. Negative
$E_0^B$ for all $N_f\ge2$ = many-magnon binding (quantum skyrmion). $\Cperp$ peaks at $0.73$
($N_f=4$), inside the paper's $0.6$--$0.8$ window.}
\label{tab:binding19}
\end{table}

\subsection{Quantitative anchor: $\Cperp = 0.73$}
The maximal transverse anticorrelation $\Cperp$ [Eq.~\eqref{eq:cperp}] rises with $N_f$ to a peak of
$\mathbf{0.73}$ at $N_f=4$ on the 19-site flake before falling as the droplet fills the flake. This
sits inside the paper's quoted $\Cperp\approx 0.6$--$0.8$ range for a genuine quantum skyrmion
obeying the transverse-correlation selection rule (Figs.~2,~5); the pdftotext extraction of Fig.~2
even preserves the tabulated correlation magnitudes $0.02,\,0.24,\,0.61,\,0.73$, confirming the
target. This is our single clean digit-level agreement with the paper.

\subsection{Scalar chirality is exactly zero --- and that is correct}
We measured the raw scalar chirality
$\langle\chi\rangle = \sum_\triangle \langle\Svec_i\cdot(\Svec_j\times\Svec_k)\rangle$ and obtained
\emph{exactly zero} in every sector. \textbf{This is correct physics, not a bug and not a miss:}
\begin{itemize}
\item $H$ is real-symmetric $\Rightarrow$ its ground state can be chosen real. The triple product
$\Svec_i\cdot(\Svec_j\times\Svec_k)$ contains an odd number of $S^y$ operators and is therefore a
purely imaginary (anti-Hermitian $\times i$) operator; its expectation in a real state is
identically zero.
\item Physically: with no spin-orbit coupling, skyrmion and antiskyrmion are \emph{exactly
degenerate} (paper Sec.~I.A), so any net chirality cancels between the two.
\end{itemize}
The paper itself does \emph{not} use raw $\langle\chi\rangle$ as the skyrmion diagnostic --- it uses
the arctan winding-correlation (Eq.~12) and the transverse antiferromagnetic correlation. We
therefore report $\langle\chi\rangle=0$ as an \emph{expected-zero}, and use $\Cperp$ as the
quantitative anchor.

\subsection{Finite-size scaling: exponentially small tunneling bandwidth (COVERAGE-FLIP)}
The paper's headline \emph{dynamical} claim (abstract) is that the quantum skyrmion's bandwidth is
``exponentially small and arises from tunneling processes between skyrmion and antiskyrmion.'' On a
finite flake with (approximate) $C_6$ symmetry, skyrmion and antiskyrmion are related by symmetry
(time reversal $\times$ spin-$\pi$ rotation, paper Sec.~I.A) and are therefore \emph{exactly
degenerate}; their residual splitting on the finite cluster is the finite-size tunneling splitting
$\Delta_{\rm tun}=E_1-E_0$ \emph{within} the skyrmion $N_f$-sector. We resolve it by computing the
lowest $k=4$--$6$ eigenpairs (not just the ground state) via \texttt{eigsh}, on flakes
$N=7,19,37$ ($37$ approaches the paper's 31-site cluster; largest sector dim $\approx 4.4\times10^5$,
$\sim46$~s). A vectorized $\Cperp$ (numpy bit-ops $+$ \texttt{searchsorted}) keeps the run inside a
$\sim74$~s budget.

Table~\ref{tab:fss} shows the result. At $N=7$ there is \emph{no} skyrmion: $\Cperp=0$ and the
level spacing is generic ($\Delta_{\rm tun}\sim0.23$). At $N=19$ and $N=37$ the skyrmion sectors
show $\Cperp$ inside the paper's $0.6$--$0.8$ window \emph{and} the low-energy spectrum collapses to
a near-degenerate skyrmion/antiskyrmion doublet with $\Delta_{\rm tun}$ falling from $\sim10^{-2}$
to \emph{machine-zero} $\sim10^{-14}$ --- an exponentially small bandwidth, exactly as claimed.
The splitting also reproduces the paper's $N_f\bmod3$ \emph{selection rule} (Figs.~7,~8): sectors
with $N_f=2\bmod3$ are exactly degenerate ($\Delta_{\rm tun}\sim10^{-14}$, quadratic/vanishing
tunneling), while $N_f=0,1\bmod3$ retain a finite $\sim10^{-2}$ splitting.

\begin{table}[h]
\centering
\begin{tabular}{rrrrl}
\toprule
$N$ & sector $N_f$ & $\Delta_{\rm tun}=E_1-E_0$ & $\Cperp$ & skyrmion? \\
\midrule
7  & 4 & $2.3\times10^{-1}$ & $0.00$ & no (generic spacing) \\
19 & 3 & $5.3\times10^{-3}$ & $0.60$ & yes ($N_f{=}0\bmod3$, finite split) \\
19 & 4 & $4.0\times10^{-14}$ & $\mathbf{0.73}$ & yes ($N_f{=}1\bmod3$) \\
19 & 5 & $3.6\times10^{-14}$ & $0.63$ & yes ($N_f{=}2\bmod3$, degenerate) \\
37 & 4 & $1.5\times10^{-2}$ & $0.40$ & yes ($N_f{=}1\bmod3$, finite split) \\
37 & 5 & $8.0\times10^{-14}$ & $1.09$ & yes ($N_f{=}2\bmod3$, degenerate) \\
\bottomrule
\end{tabular}
\caption{Finite-size scaling of the skyrmion-sector tunneling splitting
$\Delta_{\rm tun}=E_1-E_0$ ($J_1{=}1$, $J_2{=}0.5$, $K{=}0.05$, $B{=}0$). The bandwidth is
exponentially small ($\Delta_{\rm tun}$ down to $\sim10^{-14}$) once a skyrmion forms, and follows
the paper's $N_f\bmod3$ selection structure. Data: \texttt{report/evidence/lohani\_fss\_result.json}.}
\label{tab:fss}
\end{table}

\section{Comparison to the paper}
\begin{table}[h]
\centering
\begin{tabular}{p{6.2cm}p{4.2cm}l}
\toprule
Claim & Paper & This work \\
\midrule
Many-magnon bound state ($E_0^B<0$, $N_f\ge2$) & Central ED result & \textbf{Reproduced} (7 \& 19 site) \\
Transverse anticorrelation $\Cperp$ of skyrmion & $\approx 0.6$--$0.8$ (Figs.~2,5) & $\mathbf{0.73}$ ($N_f{=}4$, 19-site) \\
Raw scalar chirality & not used (Sz degeneracy) & $0$ exactly (expected) \\
Droplet forms at finite $N_f$ & $N_f^{\min}\sim5$--$9$ & $\Delta E<0$ at $N_f=5,6$ \\
FM wins global ground state at finite $B$ & yes & yes (sector sweep) \\
Exponentially small tunneling bandwidth & headline dynamical result & \textbf{Reproduced}: $\Delta_{\rm tun}\!\to\!10^{-14}$ (FSS to $N{=}37$) \\
$N_f\bmod3$ tunneling selection rule & Figs.~7,~8 & \textbf{Reproduced} (see Table~\ref{tab:fss}) \\
\bottomrule
\end{tabular}
\caption{Claim-by-claim comparison. The core ED physics and the one quantitative anchor reproduce.}
\label{tab:compare}
\end{table}

\section{What was \emph{not} reproduced (scope)}
This replication targets the ED core and its finite-size scaling. The following elements remain out
of scope and are honestly flagged as gaps (see \texttt{failure\_analysis.md}):
the exact 31-site flake with explicit $l_z$ angular-momentum symmetry labels (we reached the
\emph{larger} $N=37$ round flake and observe the $N_f\bmod3$ tunneling structure empirically, but
have not block-diagonalized by an exact $C_6$ operator); the full $J_2$--$B$ phase diagram; the
phenomenological Schr\"odinger equation for skyrmion motion; and the Eq.~12 arctan winding-number
form. These require symmetry-adapted bases, 2D parameter sweeps, or a separate effective-model build.

\section{Verdict}
\textbf{REPLICATED} (core exact-diagonalization claim, the quantitative $\Cperp$ anchor, \emph{and}
the headline dynamical claim --- exponentially small skyrmion--antiskyrmion tunneling bandwidth ---
via larger-flake finite-size scaling).
\begin{itemize}
\item \textbf{Coverage} $\approx 7/10$: Hamiltonian, geometry, binding-energy criterion,
transverse-correlation diagnostic, larger flake ($N=37$), and the tunneling-splitting FSS all built
and verified; the exact-$C_6$ $l_z$ labels, full phase diagram, and Eq.~12 winding form not built.
\item \textbf{Agreement} $\approx 8/10$: $\Cperp=0.73$ lands inside the paper's $0.6$--$0.8$ window;
$E_0^B<0$ for all $N_f\ge2$ and deepens with $N_f$; droplet excitation energy goes negative at
$N_f=5,6$; and the skyrmion-sector tunneling splitting collapses to $\sim10^{-14}$ with the correct
$N_f\bmod3$ selection structure. Not a digit-for-digit match because the flake size/geometry differ
from the paper's exact 31-site cluster.
\end{itemize}

\section*{Reproducibility}
Code: \texttt{report/evidence/lohani\_ed.py} (base ED) and
\texttt{report/evidence/lohani\_fss.py} (larger-flake FSS + tunneling splitting). Numeric output:
\texttt{report/evidence/lohani2019\_result.json} (includes the
\texttt{finite\_size\_tunneling\_splitting} section) and
\texttt{report/evidence/lohani\_fss\_result.json}. Interpreter:
\texttt{/home/stevens/comfyui-env/bin/python} (\texttt{scipy.sparse.linalg.eigsh}). Base run
(7-site sweep $+$ binding analyses $+$ 19-site up to $N_f=6$) completes in $\sim9$~s; the FSS run
($N=7,19,37$) in $\sim74$~s.

\end{document}
